% --- Encoding / typography ---
\usepackage[utf8]{inputenc}
\usepackage[T1]{fontenc}
\usepackage{microtype}
\usepackage{mathptmx} % Times New Roman-like font

% --- Language ---
\usepackage[spanish]{babel}
\usepackage{csquotes}

% --- Page geometry ---
% headheight evita el warning de fancyhdr
\usepackage[top=1in, bottom=1in, left=1.5in, right=1in, headheight=16pt]{geometry}

% --- Index ---
\usepackage{makeidx}

% --- Code listings ---
\usepackage{fvextra}

\usepackage[T1]{fontenc}
\usepackage[utf8]{inputenc}
\usepackage{listings}
\usepackage{xcolor}

\lstdefinestyle{pythoncompact}{
    language=Python,
    basicstyle=\ttfamily\scriptsize,
    keywordstyle=\color{blue!70!black}\bfseries,
    commentstyle=\color{green!40!black},
    stringstyle=\color{red!60!black},
    showstringspaces=false,
    breaklines=true,
    breakatwhitespace=false,
    tabsize=4,
    frame=single,
    rulecolor=\color{black!30},
    captionpos=b,
    numbers=none,
    xleftmargin=0pt,
    framexleftmargin=0pt,
    aboveskip=4pt,
    belowskip=4pt,
    lineskip=-1pt,
    columns=fullflexible,
    inputencoding=utf8,
    extendedchars=true,
    literate=
        {á}{{\'a}}1
        {é}{{\'e}}1
        {í}{{\'i}}1
        {ó}{{\'o}}1
        {ú}{{\'u}}1
        {Á}{{\'A}}1
        {É}{{\'E}}1
        {Í}{{\'I}}1
        {Ó}{{\'O}}1
        {Ú}{{\'U}}1
        {ñ}{{\~n}}1
        {Ñ}{{\~N}}1
}

\renewcommand{\lstlistingname}{Listado}
\renewcommand{\lstlistlistingname}{Lista de listados}

% --- Tables ---
\usepackage{tabularx}
\newcolumntype{B}{>{\hsize=1.5\hsize}X}
\newcolumntype{s}{>{\hsize=.5\hsize}X}
\usepackage{booktabs}
\usepackage{float}
\usepackage{multirow}
\usepackage{longtable}
% --- Bibliography ---
\usepackage[backend=biber, style=ieee]{biblatex}
\addbibresource{Bibliografia.bib}
\usepackage{xurl}
\usepackage{hyperref}

% --- Chapter formatting ---
\usepackage{titlesec}
\titleformat{\chapter}[display]
  {\normalfont\Large\bfseries}
  {\chaptertitlename\ \thechapter}
  {20pt}
  {\LARGE}
\titlespacing*{\chapter}{0pt}{-30pt}{40pt}

% --- TOC formatting ---
\usepackage[titles]{tocloft}
\renewcommand{\cftchapfont}{\normalfont\bfseries}
\renewcommand{\cftchapleader}{\cftdotfill{\cftdotsep}}
\addto\captionsspanish{\renewcommand{\tablename}{Tabla}}

% --- Figures / captions ---
\usepackage{graphicx}
\graphicspath{{./figuras/}}
\usepackage[font=small, labelfont=bf]{caption}

% --- Math ---
\usepackage{amsmath}
\usepackage{amssymb}
\usepackage{amsthm}

% --- Theorems and definitions ---
\newtheorem{theorem}{Theorem}[chapter]
\newtheorem{lemma}[theorem]{Lemma}
\newtheorem{proposition}[theorem]{Proposition}
\newtheorem{corollary}[theorem]{Corollary}
\theoremstyle{definition}
\newtheorem{definition}[theorem]{Definition}
\newtheorem{example}[theorem]{Example}

% --- Header and footer (fancyhdr) ---
\usepackage{fancyhdr}
\pagestyle{fancy}
\fancyhf{}

% En report, \leftmark suele ser "Capítulo N. Título". Esto lo deja consistente.
\renewcommand{\chaptermark}[1]{\markboth{\thechapter.\ #1}{}}

\fancyhead[L]{\nouppercase{\leftmark}}
\fancyfoot[C]{\thepage}
\renewcommand{\headrulewidth}{0.4pt}
\renewcommand{\footrulewidth}{0pt}

% --- Line spacing ---
\usepackage{setspace}
\onehalfspacing



\newcommand{\imagen}[2]{%
    \includegraphics[
        width=#1\textwidth,
        height=0.82\textheight,
        keepaspectratio
    ]{#2}%
}

\newcommand{\p}[1]{%
    \par\vspace{#1pt}%
}

\newcommand{\figura}[5]{%
    \begin{figure}[H]
        \centering
        \imagen{#1}{#2}\p{5}
        \caption[#3]{#3 #4}
        #5
    \end{figure}
}